\section{Discussion}
\label{sec:discussion}
\subsection{What the matched comparisons establish}
The seed44 analysis shows a descriptive increase in mean AUROC and AP as the model progresses from direct prediction without future-state supervision to auxiliary-supervised direct prediction and then to forecast-attention readout. The first contrast is consistent with future-state supervision encouraging an observed-history representation that is useful for pCR prediction. The second is consistent with additional predictive value in a classifier that receives the forecast state and its change. However, the latter comparison also changes classifier width and pooling, so it cannot attribute the difference exclusively to a more accurate latent forecast. The supplementary mean-pooling control narrows the pooling question but exhibits only a small mean AUROC difference.

This distinction also changes how the model should be described. The main experiments do not reconstruct future MRI for classification and do not establish that an imagined image reproduces a patient's future anatomy. Rather, an auxiliary latent objective regularizes a representation, and the full classifier reads a predicted embedding. Moreover, the direct control named \nojepa{} retains the common fusion backbone and its transition queries. It is an ablation of future-state supervision, not proof that every world-model-related architectural component has been removed. Describing the actual information paths is more informative than relying on the experiment labels alone.

\subsection{Longitudinal prediction and the role of treatment context}
The full model's AUROC increased across T0--T2 in the selected seed, and the larger differences relative to the direct control occurred at later landmarks. One plausible explanation is that an observed response trajectory helps condition later-state forecasts. Another is that the additional forecast features or classifier capacity suit the later-prefix distribution. The current experiments do not distinguish these explanations. Forecast-quality analysis, frozen-forecast controls, parameter-matched direct classifiers, and prospective landmark-specific validation would help determine the mechanism.

Treatment is represented as an assigned-arm scenario assumed known at baseline. In the source protocol, T0 MRI preceded randomization; therefore, using the eventual assignment with T0 MRI is a retrospective conditional analysis and cannot be interpreted as a deployable prediction before allocation. The input does not encode the patient's verified delivered dose, dose reductions, actual cycle timing, adherence, or later treatment changes. Consequently, the framework should be interpreted as response prediction conditional on an encoded assignment, not as a simulator for treatment intervention. Comparing predictions under different assigned-arm tokens would not estimate causal treatment effects without additional assumptions, appropriate treatment data, and a separate causal analysis.

\subsection{Discrimination is not calibrated clinical utility}
The forecast-attention model's improved AUROC and AP coexisted with lower mean F1 than auxiliary-supervised direct readout and worse NLL, Brier score, and ECE than the control without future-state supervision. These discrepancies are expected to matter if probabilities are used to support decisions. AUROC measures ranking over thresholds, while F1 at 0.5 depends on a particular operating point; neither alone certifies useful clinical risk estimates. The model was selected by validation F1, not by a prespecified utility function. Independent assessment of calibration, clinically motivated operating points, and decision consequences is needed before suggesting a role in patient management. The attention maps also require separate validation before being interpreted as anatomical explanations.

\subsection{Limitations and requirements for a confirmatory study}
\label{sec:limitations}
The most immediate limitation is retrospective seed selection. Seed44 was chosen for manuscript presentation after multiple-seed results were available, including holdout results. A matched seed does not make this selection prospective. Patient bootstrapping quantifies variation in the evaluated patient sample conditional on the trained models, but cannot account for selection among seeds or the randomness of retraining. All reported model comparisons and intervals should therefore be considered exploratory. A confirmatory study should lock the model, seed policy, endpoints, and analysis before evaluating new patients.

The second limitation concerns the evaluation cohort and upstream independence. The holdout is an internal downstream split, not an external institution or independent prospective cohort. Twenty-nine of its patients had belonged to a development split in an earlier experimental lineage. Current supervised training uses the new patient-disjoint folds and does not warm-start from those earlier pCR models. Nevertheless, the patient-level overlap of the pretrained codec and localization-segmentation development data with the current holdout has not been fully established. This unresolved provenance prevents a claim that the complete pipeline has been assessed on entirely unseen patients at every stage.

Third, the fixed crop and preprocessing introduce constraints. A T0-derived localization window is used for all visits, and some predicted lesion components extend beyond the fixed field of view. Geometry and array-consistency checks do not substitute for anatomical review of segmentation or registration. The spatial window may underrepresent multifocal disease, chest-wall extent, or changes outside the crop. Complete T0--T2 availability is also required for holdout inclusion, whereas the development set retains incomplete histories. Performance may therefore differ for patients with missing interim MRI, and the missingness process may itself be associated with clinical factors.

Fourth, the present empirical comparison comprises matched internal variants. It does not yet establish performance against clinical-only models, conventional imaging biomarkers, independent longitudinal architectures, or contemporary published methods under the same split. Patient-subgroup analysis, scanner and site stratification, subtype-specific sensitivity, and the relationship between forecast error and classification error are also incomplete. The relatively small holdout, with 56 positive outcomes, limits the precision of those analyses and the principal paired estimates.

Finally, several reporting details still require confirmation from original clinical documentation: the exact pCR pathology definition and assessment workflow, actual patient-level acquisition intervals and treatment exposure, clinical-data adjudication, and the applicable ethics determination. The nominal stage timing is documented, but adherence and interval distributions have not yet been audited. These are explicit outstanding items in this draft rather than assumptions filled from common practice. Independent external validation, complete pretraining provenance, suitable external baselines, and a prespecified analysis would materially strengthen the evidence for a final Medical Image Analysis submission.
